% !TEX spellcheck = en_US



% ~~~~~~~~~~~~~~~~~~~~~~~~~~~~~~~~~~~~~~~ V
% (NC) 2026 Haluk Bingol
% github.com/halukbingol/LaTeX-Templates
%
% Licensees may copy, distribute, display, and perform the work and make derivative works 
% and remixes based on it only for non-commercial purposes. 
% ~~~~~~~~~~~~~~~~~~~~~~~~~~~~~~~~~~~~~~~ A




% =====================================================================
%  Internationalization in Java -- Junior Level: Unicode, Locale Data and Time in Depth
%  ---------------------------------------------------------------------
%  Built on Template/hblibTemplate.tex with the libraries in ../../hblib/.
%  Programs, terminal sessions and their real outputs are in codeoutput/:
%     codeoutput/<Name>.java  a program        codeoutput/<Name>.sh   a session
%     codeoutput/<Name>/      the session's project folder
%     codeoutput/<Name>.txt   what it printed, made by:  bash run_examples.sh
% =====================================================================

\documentclass[aspectratio=169,10pt,compress]{beamer}

% ~~~~~~~~~~~~~~~~~~~~~~~~~~~~~~~~~~~~~~~~~~~~~~~~~~~~~~~~~~~~~~~~~ V hblib
	\usepackage{../../hblib/hblibPresentation} % lib presentation: theme, i/N, \hSection, algorithm2e, ...
	\usepackage{../../hblib/hblibCodeoutput} % lib code + output: \lstcode, \lstterminal, ...
	\usepackage{../../hblib/hblibDatastructures} % lib data structures: \hString, \hStack, \hQueue
% ~~~~~~~~~~~~~~~~~~~~~~~~~~~~~~~~~~~~~~~~~~~~~~~~~~~~~~~~~~~~~~~~~ A hblib


% xcolor, array (and the L{width} column), listings, tikz and algorithm2e come from the hblib libraries
\usepackage[T1]{fontenc}
\usepackage{lmodern}
\usepackage{booktabs}

\definecolor{gitred}{RGB}{240,80,50}
\newcommand{\win}{\textbf{Windows}}
\newcommand{\mac}{\textbf{macOS}}

% algorithm2e set-up (the package is loaded by hblibPresentation), as in the template
\SetAlCapFnt{\small}
\SetAlCapNameFnt{\small}
\DontPrintSemicolon

% .properties files (resource bundles)
\lstdefinelanguage{properties}{morecomment=[l]{\#}, morecomment=[l]{!}}

% ---------------------------------------------------------------------
\title{Internationalization in Java}
\subtitle{Junior Level: Unicode, Locale Data and Time in Depth}
\author[Bingol]{Haluk O. Bingol}

\institute[FCIS]{
    Faculty of Computer and Information Sciences\\
    Yeditepe University
}
\date{
	{\tiny \hbTimeStamp}
}

\begin{document}




% ~~~~~~~~~~~~~~~~~~~~~~~~~~~~~~~~~~~~~~ V frm
\begin{frame}[t,fragile]
  \titlepage
\end{frame}
% ~~~~~~~~~~~~~~~~~~~~~~~~~~~~~~~~~~~~~~ A frm




% ~~~~~~~~~~~~~~~~~~~~~~~~~~~~~~~~~~~~~~ V frm
\begin{frame}[t,fragile] \frametitle{Outline}
  \begin{columns}[T]
    \begin{column}{0.48\textwidth}
      \tableofcontents[sections={1-3}, hideallsubsections]
    \end{column}
    \begin{column}{0.48\textwidth}
      \tableofcontents[sections={4-20}, hideallsubsections]
    \end{column}
  \end{columns}
\end{frame}
% ~~~~~~~~~~~~~~~~~~~~~~~~~~~~~~~~~~~~~~ A frm

% ~~~~~~~~~~~~~~~~~~~~~~~~~~~~~~~~~~~~~~ V
\hSection[Unicode]{Unicode in Depth}
% ~~~~~~~~~~~~~~~~~~~~~~~~~~~~~~~~~~~~~~ A

% ~~~~~~~~~~~~~~~~~~~~~~~~~~~~~~~~~~~~~~ V
\subsection[Units]{Bytes, chars, code points, graphemes}
% ~~~~~~~~~~~~~~~~~~~~~~~~~~~~~~~~~~~~~~ A




% ~~~~~~~~~~~~~~~~~~~~~~~~~~~~~~~~~~~~~~ V frm
\begin{frame}[t,fragile] \frametitle{Four ways to count a string}
  \begin{columns}[T]
    \begin{column}{0.54\textwidth}
      \begin{center}\scriptsize
      \begin{tabular}{ll}
        \toprule
        \textbf{unit} & \textbf{Java} \\
        \midrule
        byte       & \texttt{getBytes(UTF\_8).length} \\
        \texttt{char} (UTF-16 unit) & \texttt{length()}, \texttt{charAt(i)} \\
        code point & \texttt{codePointCount}, \texttt{codePoints()} \\
        grapheme   & \texttt{BreakIterator.getCharacterInstance()} \\
        \bottomrule
      \end{tabular}
      \end{center}
      \small A \textbf{grapheme} is what a reader sees as one character: \texttt{é} written as \texttt{e} + accent,
      or a thumbs-up with a skin tone (2 code points, 4 \texttt{char}s).
    \end{column}
    \begin{column}{0.42\textwidth}
      \begin{itemize}\small
        \item Code points above U+FFFF (most emoji, historic scripts, rare CJK) need \textbf{two} \texttt{char}s:
          a \emph{surrogate pair}.
        \item \texttt{substring}, \texttt{charAt}, \texttt{reverse} can cut a character in half.
        \item Cursor movement, ``max 20 characters'', truncation with \dots{} should count \textbf{graphemes}.
        \item Since Java 20, \texttt{BreakIterator} follows the Unicode rules for emoji sequences.
      \end{itemize}
    \end{column}
  \end{columns}
\end{frame}
% ~~~~~~~~~~~~~~~~~~~~~~~~~~~~~~~~~~~~~~ A frm




% ~~~~~~~~~~~~~~~~~~~~~~~~~~~~~~~~~~~~~~ V frm
\begin{frame}[t,fragile] \frametitle{Counting in practice}
  \lstcode[basicstyle=\ttfamily\tiny, firstline=9, firstnumber=9]{codeoutput/Units.java}{Java}
  \uncover<2->{\lstterminal[basicstyle=\ttfamily\tiny]{codeoutput/Units.txt}}

  \vspace{0.1em}
  \small
  Thumbs-up + skin tone: 6 \texttt{char}s, 4 code points, 3 graphemes.
\end{frame}
% ~~~~~~~~~~~~~~~~~~~~~~~~~~~~~~~~~~~~~~ A frm

% ~~~~~~~~~~~~~~~~~~~~~~~~~~~~~~~~~~~~~~ V
\subsection[Normalization]{Normalization forms}
% ~~~~~~~~~~~~~~~~~~~~~~~~~~~~~~~~~~~~~~ A




% ~~~~~~~~~~~~~~~~~~~~~~~~~~~~~~~~~~~~~~ V frm
\begin{frame}[t,fragile] \frametitle{Four normalization forms}
  \begin{columns}[T]
    \begin{column}{0.5\textwidth}
      \lstcode[basicstyle=\ttfamily\tiny, firstline=4, firstnumber=4]{codeoutput/Forms.java}{Java}
    \end{column}
    \begin{column}{0.48\textwidth}
      \uncover<2->{\lstterminal[basicstyle=\ttfamily\tiny]{codeoutput/Forms.txt}}
    \end{column}
  \end{columns}

  \vspace{0.1em}
  \small \textbf{NFC} composes, \textbf{NFD} decomposes. The \textbf{K} forms also replace
  \emph{compatibility} characters: ligature \texttt{FB01}\,$\to$\,\texttt{fi}, full-width \texttt{FF21 FF22 FF11}\,$\to$\,\texttt{AB1},
  \texttt{x²}\,$\to$\,\texttt{x2}. Store NFC; compare identifiers (user names) after NFKC.
\end{frame}
% ~~~~~~~~~~~~~~~~~~~~~~~~~~~~~~~~~~~~~~ A frm

% ~~~~~~~~~~~~~~~~~~~~~~~~~~~~~~~~~~~~~~ V
\subsection[Case]{Case mapping}
% ~~~~~~~~~~~~~~~~~~~~~~~~~~~~~~~~~~~~~~ A




% ~~~~~~~~~~~~~~~~~~~~~~~~~~~~~~~~~~~~~~ V frm
\begin{frame}[t,fragile] \frametitle{Case mapping is not one-to-one}
  \lstcode[basicstyle=\ttfamily\tiny, firstline=10, firstnumber=10]{codeoutput/Mapping.java}{Java}
  \uncover<2->{\lstterminal[basicstyle=\ttfamily\tiny]{codeoutput/Mapping.txt}}

  \vspace{0.1em}
  \small
  The length can change (\texttt{ß}\,$\to$\,\texttt{SS}); \texttt{İ} lower-cased with \texttt{ROOT} is \texttt{i} + U+0307;
  \texttt{equalsIgnoreCase} works char by char and misses \texttt{ß}/\texttt{SS}.
\end{frame}
% ~~~~~~~~~~~~~~~~~~~~~~~~~~~~~~~~~~~~~~ A frm

% ~~~~~~~~~~~~~~~~~~~~~~~~~~~~~~~~~~~~~~ V
\hSection[Locale data]{Locale Data in Depth}
% ~~~~~~~~~~~~~~~~~~~~~~~~~~~~~~~~~~~~~~ A

% ~~~~~~~~~~~~~~~~~~~~~~~~~~~~~~~~~~~~~~ V
\subsection[Extensions]{Unicode extensions}
% ~~~~~~~~~~~~~~~~~~~~~~~~~~~~~~~~~~~~~~ A




% ~~~~~~~~~~~~~~~~~~~~~~~~~~~~~~~~~~~~~~ V frm
\begin{frame}[t,fragile] \frametitle{Unicode locale extensions: \texttt{-u-}}
  \lstcode[basicstyle=\ttfamily\tiny, firstline=8, firstnumber=8]{codeoutput/Extensions.java}{Java}
  \uncover<2->{\lstterminal[basicstyle=\ttfamily\tiny]{codeoutput/Extensions.txt}}

  \vspace{0.1em}
  \small
  \texttt{ca} calendar, \texttt{cu} currency, \texttt{fw} first day, \texttt{rg} region; \texttt{localizedBy} (Java 10) applies them.
\end{frame}
% ~~~~~~~~~~~~~~~~~~~~~~~~~~~~~~~~~~~~~~ A frm

% ~~~~~~~~~~~~~~~~~~~~~~~~~~~~~~~~~~~~~~ V
\subsection[Categories]{Display and format}
% ~~~~~~~~~~~~~~~~~~~~~~~~~~~~~~~~~~~~~~ A




% ~~~~~~~~~~~~~~~~~~~~~~~~~~~~~~~~~~~~~~ V frm
\begin{frame}[t,fragile] \frametitle{Two default locales}
  \begin{columns}[T]
    \begin{column}{0.5\textwidth}
      \lstcode[basicstyle=\ttfamily\tiny]{codeoutput/Categories/Cat.java}{Java}
    \end{column}
    \begin{column}{0.48\textwidth}
      \uncover<2->{\lstterminal[basicstyle=\ttfamily\tiny]{codeoutput/Categories.txt}}
    \end{column}
  \end{columns}

  \vspace{0.1em}
  \small \texttt{Locale.Category.DISPLAY} names things (languages, months in menus),
  \texttt{FORMAT} formats values. A user may want English menus with Turkish numbers:
  \texttt{-Duser.language.format=tr} (or \texttt{Locale.setDefault(Category.FORMAT, l)}).
\end{frame}
% ~~~~~~~~~~~~~~~~~~~~~~~~~~~~~~~~~~~~~~ A frm

% ~~~~~~~~~~~~~~~~~~~~~~~~~~~~~~~~~~~~~~ V
\subsection[Providers]{Where the data comes from}
% ~~~~~~~~~~~~~~~~~~~~~~~~~~~~~~~~~~~~~~ A




% ~~~~~~~~~~~~~~~~~~~~~~~~~~~~~~~~~~~~~~ V frm
\begin{frame}[t,fragile] \frametitle{Locale data providers: CLDR and COMPAT}
  \begin{columns}[T]
    \begin{column}{0.48\textwidth}
      \lstcode[basicstyle=\ttfamily\tiny]{codeoutput/Providers/Prov.java}{Java}
    \end{column}
    \begin{column}{0.50\textwidth}
      \uncover<2->{\lstterminal[basicstyle=\ttfamily\tiny]{codeoutput/Providers.txt}}
    \end{column}
  \end{columns}

  \vspace{0.1em}
  \small Since Java 9 the data comes from \textbf{CLDR}. The old JRE data (\texttt{COMPAT}) is deprecated
  in Java 21 and removed in Java 23. Upgrading Java can change formatted output: never parse your own UI text in tests.
\end{frame}
% ~~~~~~~~~~~~~~~~~~~~~~~~~~~~~~~~~~~~~~ A frm

% ~~~~~~~~~~~~~~~~~~~~~~~~~~~~~~~~~~~~~~ V
\hSection[Plurals]{Plural Rules}
% ~~~~~~~~~~~~~~~~~~~~~~~~~~~~~~~~~~~~~~ A




% ~~~~~~~~~~~~~~~~~~~~~~~~~~~~~~~~~~~~~~ V frm
\begin{frame}[t,fragile] \frametitle{CLDR plural rules in the JDK}
  \lstcode[basicstyle=\ttfamily\tiny, firstline=4, firstnumber=4]{codeoutput/Plurals.java}{Java}
  \uncover<2->{\lstterminal[basicstyle=\ttfamily\tiny]{codeoutput/Plurals.txt}}

  \vspace{0.1em}
  \small
  English and German: 2 forms; Turkish uses one; Polish has 3 (\emph{milion}, \emph{miliony}, \emph{milionów}:
  1 / 2--4, 22 / 5, 25). \texttt{ChoiceFormat} cannot express Polish or Arabic rules; use ICU4J or keep messages simple.
\end{frame}
% ~~~~~~~~~~~~~~~~~~~~~~~~~~~~~~~~~~~~~~ A frm




% ~~~~~~~~~~~~~~~~~~~~~~~~~~~~~~~~~~~~~~ V frm
\begin{frame}[t,fragile] \frametitle{Messages with real plural rules}
  \begin{columns}[T]
    \begin{column}{0.54\textwidth}


% +++++++++++++++++++++++++++++++++++++++ V lst
%: -Lst.
\begin{lstlisting}[style=codesnippet, language=Java, basicstyle=\ttfamily\scriptsize]
// ICU4J (com.ibm.icu:icu4j), not in the JDK
MessageFormat m = new MessageFormat(
    "{n, plural, one {# dosya} other {# dosya}}",
    ULocale.forLanguageTag("tr"));
// Polish: one / few / many / other
"{n, plural, one {# plik} few {# pliki} many {# plików} other {# pliku}}"
\end{lstlisting}
% +++++++++++++++++++++++++++++++++++++++ A lst


      \small \hRemark{Not run here} (ICU4J comes from Maven Central).
    \end{column}
    \begin{column}{0.42\textwidth}
      \begin{itemize}\small
        \item CLDR categories: \texttt{zero}, \texttt{one}, \texttt{two}, \texttt{few}, \texttt{many}, \texttt{other}.
        \item Each language uses some of them; the translator fills in the forms.
        \item The JDK has the rules (used by \texttt{CompactNumberFormat}) but no public plural API;
          \texttt{java.text.ListFormat} (\emph{a, b and c}) arrived in Java 22.
      \end{itemize}
    \end{column}
  \end{columns}
\end{frame}
% ~~~~~~~~~~~~~~~~~~~~~~~~~~~~~~~~~~~~~~ A frm

% ~~~~~~~~~~~~~~~~~~~~~~~~~~~~~~~~~~~~~~ V
\hSection[Lookup]{Bundle Lookup in Detail}
% ~~~~~~~~~~~~~~~~~~~~~~~~~~~~~~~~~~~~~~ A




% ~~~~~~~~~~~~~~~~~~~~~~~~~~~~~~~~~~~~~~ V frm
\begin{frame}[t,fragile] \frametitle{Candidate locales}
  \begin{columns}[T]
    \begin{column}{0.5\textwidth}
      \lstcode[basicstyle=\ttfamily\tiny]{codeoutput/Candidates.java}{Java}
    \end{column}
    \begin{column}{0.48\textwidth}
      \uncover<2->{\lstterminal[basicstyle=\ttfamily\tiny]{codeoutput/Candidates.txt}}
    \end{column}
  \end{columns}

  \vspace{0.1em}
  \small Script and variant are tried first; special cases: \texttt{zh-TW} implies the script \texttt{Hant},
  Norwegian Bokmål \texttt{nb} also tries the old code \texttt{no}. The empty locale is the base bundle.
\end{frame}
% ~~~~~~~~~~~~~~~~~~~~~~~~~~~~~~~~~~~~~~ A frm




% ~~~~~~~~~~~~~~~~~~~~~~~~~~~~~~~~~~~~~~ V frm
\begin{frame}[t,fragile] \frametitle{\texttt{ResourceBundle.getBundle}}
% +++++++++++++++++++++++++++++++++++++++ V alg
%: -Alg. 
\begin{algorithm}[H]
  \caption{\textsc{GetBundle}$(b, \ell)$ (default \texttt{Control})}
  \KwIn{base name $b$, locale $\ell$, default locale $\delta$, class loader}
  \For{$t \gets \ell, \delta$}{
    \ForEach{$c$ in \textsc{CandidateLocales}$(b, t)$}{
      \tcp{e.g.\ \texttt{tr\_TR}, \texttt{tr}, the empty locale}
      \lIf{the cache has bundle $(b, c)$}{use it}
      \lElse{try class \texttt{b\_c}, then file \texttt{b\_c.properties}; cache the result (also ``not found'')}
      link the found bundles: each one's parent is the next found candidate\;
    }
    \lIf{a bundle other than the base was found for $t$}{\Return{the first one}}
    \lIf{$t = \delta$ or \textsc{FallbackLocale}$(b, \ell) = $ none}{\textbf{break}}
  }
  \Return{the base bundle, or throw \texttt{MissingResourceException}}\;
\end{algorithm}
% +++++++++++++++++++++++++++++++++++++++ A alg

  \small Bundles are \textbf{cached}; \texttt{ResourceBundle.clearCache()} or a \texttt{Control} with a time-to-live reloads them.
\end{frame}
% ~~~~~~~~~~~~~~~~~~~~~~~~~~~~~~~~~~~~~~ A frm

% ~~~~~~~~~~~~~~~~~~~~~~~~~~~~~~~~~~~~~~ V
\hSection[Time]{Time Zones in Depth}
% ~~~~~~~~~~~~~~~~~~~~~~~~~~~~~~~~~~~~~~ A




% ~~~~~~~~~~~~~~~~~~~~~~~~~~~~~~~~~~~~~~ V frm
\begin{frame}[t,fragile] \frametitle{Time zone rules change}
  \lstcode[basicstyle=\ttfamily\tiny, firstline=8, firstnumber=8]{codeoutput/TimeRules.java}{Java}
  \uncover<2->{\lstterminal[basicstyle=\ttfamily\tiny]{codeoutput/TimeRules.txt}}

  \vspace{0.1em}
  \small
  Zone rules are law (Türkiye: no clock changes since 2016). Autumn: \texttt{02:30} happens twice.
\end{frame}
% ~~~~~~~~~~~~~~~~~~~~~~~~~~~~~~~~~~~~~~ A frm




% ~~~~~~~~~~~~~~~~~~~~~~~~~~~~~~~~~~~~~~ V frm
\begin{frame}[t,fragile] \frametitle{Which type for what?}
  \begin{center}\small
  \begin{tabular}{L{0.22\textwidth}L{0.38\textwidth}L{0.3\textwidth}}
    \toprule
    \textbf{type} & \textbf{meaning} & \textbf{use for} \\
    \midrule
    \texttt{Instant}        & a point on the time line (UTC)        & logs, database, events \\
    \texttt{ZonedDateTime}  & instant + zone rules                  & meetings, display, arithmetic over DST \\
    \texttt{OffsetDateTime} & instant + fixed offset                & data exchange (ISO 8601) \\
    \texttt{LocalDateTime}  & wall clock, no zone                   & ``every day at 09:00'', forms \\
    \texttt{LocalDate}      & a day                                 & birthdays, holidays \\
    \bottomrule
  \end{tabular}
  \end{center}
  \begin{itemize}\small
    \item Store \texttt{Instant}s; keep the user's \texttt{ZoneId} as a preference; convert only to show.
    \item For a future appointment store the local time \emph{and} the zone: the law may change before the date.
  \end{itemize}
\end{frame}
% ~~~~~~~~~~~~~~~~~~~~~~~~~~~~~~~~~~~~~~ A frm

% ~~~~~~~~~~~~~~~~~~~~~~~~~~~~~~~~~~~~~~ V
\hSection[Threads]{Formatters and Threads}
% ~~~~~~~~~~~~~~~~~~~~~~~~~~~~~~~~~~~~~~ A




% ~~~~~~~~~~~~~~~~~~~~~~~~~~~~~~~~~~~~~~ V frm
\begin{frame}[t,fragile] \frametitle{Formatters and threads}
  \lstcode[basicstyle=\ttfamily\tiny, firstline=8, firstnumber=8]{codeoutput/Unsafe.java}{Java}
  \uncover<2->{\lstterminal[basicstyle=\ttfamily\tiny]{codeoutput/Unsafe.txt}}

  \vspace{0.1em}
  \small
  \texttt{SimpleDateFormat} keeps state: shared between threads it mixes results. \texttt{DateTimeFormatter} is immutable.
\end{frame}
% ~~~~~~~~~~~~~~~~~~~~~~~~~~~~~~~~~~~~~~ A frm




% ~~~~~~~~~~~~~~~~~~~~~~~~~~~~~~~~~~~~~~ V frm
\begin{frame}[t,fragile] \frametitle{Rules for formatters}
  \begin{center}\small
  \begin{tabular}{lll}
    \toprule
    \textbf{class} & \textbf{thread-safe?} & \textbf{how to share} \\
    \midrule
    \texttt{DateTimeFormatter}           & yes  & \texttt{static final} \\
    \texttt{Locale}, \texttt{Currency}, \texttt{ZoneId} & yes & anywhere \\
    \texttt{ResourceBundle}               & yes (read only) & cache per locale \\
    \texttt{NumberFormat}, \texttt{DecimalFormat} & no & new per use, or \texttt{ThreadLocal} \\
    \texttt{SimpleDateFormat}, \texttt{MessageFormat} & no & replace by \texttt{java.time}, new per use \\
    \texttt{Collator}                     & no  & \texttt{clone()} per thread; \texttt{CollationKey} for big sorts \\
    \bottomrule
  \end{tabular}
  \end{center}
  \small \texttt{getInstance(l)} methods return a \emph{new} object each time, built from cached data: cheap enough for most code.
\end{frame}
% ~~~~~~~~~~~~~~~~~~~~~~~~~~~~~~~~~~~~~~ A frm

% ~~~~~~~~~~~~~~~~~~~~~~~~~~~~~~~~~~~~~~ V
\hSection[Direction]{Text Direction}
% ~~~~~~~~~~~~~~~~~~~~~~~~~~~~~~~~~~~~~~ A




% ~~~~~~~~~~~~~~~~~~~~~~~~~~~~~~~~~~~~~~ V frm
\begin{frame}[t,fragile] \frametitle{Right-to-left languages}
  \begin{columns}[T]
    \begin{column}{0.52\textwidth}
      \begin{itemize}\small
        \item Arabic, Hebrew, Persian, Urdu are written \textbf{right to left}; numbers and Latin words inside
          them still run left to right: \emph{bidirectional} text.
        \item The \textbf{Unicode Bidi Algorithm} decides the display order; the stored order is the reading order.
        \item \texttt{java.text.Bidi} computes the levels; Swing uses
          \texttt{ComponentOrientation.getOrientation(locale)} to mirror whole windows.
      \end{itemize}
    \end{column}
    \begin{column}{0.44\textwidth}
      \begin{block}{Design for it}
        \small
        \begin{itemize}
          \item use \emph{start}/\emph{end}, not \emph{left}/\emph{right}
          \item mirror arrows and progress bars
          \item never build a sentence from fragments: the order changes
          \item test with a right-to-left locale (\texttt{ar}, \texttt{he})
        \end{itemize}
      \end{block}
    \end{column}
  \end{columns}
\end{frame}
% ~~~~~~~~~~~~~~~~~~~~~~~~~~~~~~~~~~~~~~ A frm

% ~~~~~~~~~~~~~~~~~~~~~~~~~~~~~~~~~~~~~~ V
\hSection[Summary]{Summary}
% ~~~~~~~~~~~~~~~~~~~~~~~~~~~~~~~~~~~~~~ A




% ~~~~~~~~~~~~~~~~~~~~~~~~~~~~~~~~~~~~~~ V frm
\begin{frame}[t,fragile] \frametitle{Big ideas}
  \begin{itemize}\small
    \item A string has bytes, \texttt{char}s, code points and graphemes: count what you mean.
    \item Normalize (NFC to store, NFKC to compare identifiers) before comparing text from outside.
    \item Case mapping may change the length and depends on the language.
    \item Locales carry extensions (calendar, currency, first day, numbering); \texttt{localizedBy} uses them.
    \item Locale data comes from CLDR and changes between Java versions; time zone rules change by law.
    \item Plural rules differ widely; the JDK has them only inside \texttt{CompactNumberFormat}.
    \item Know which formatters are thread-safe; prefer \texttt{java.time}.
  \end{itemize}
\end{frame}
% ~~~~~~~~~~~~~~~~~~~~~~~~~~~~~~~~~~~~~~ A frm




% ~~~~~~~~~~~~~~~~~~~~~~~~~~~~~~~~~~~~~~ V frm
\begin{frame}[t,fragile] \frametitle{Exercises}
  \begin{enumerate}\small
    \item Write \texttt{truncate(s, n)} that keeps at most \texttt{n} graphemes and never cuts a surrogate pair.
    \item Check whether two user names are ``the same'' (NFKC + case folding). Which pairs surprise you?
    \item Print next month's calendar with the right first day of the week for \texttt{en-US}, \texttt{tr-TR}, \texttt{ar-SA}.
    \item Find a date and zone where \texttt{LocalDateTime} + one day $\neq$ \texttt{ZonedDateTime} + one day.
    \item Make \texttt{Unsafe} correct with a \texttt{ThreadLocal<SimpleDateFormat>}; measure the cost.
    \item Run your sophomore project with \texttt{-Djava.locale.providers=COMPAT}: what changes?
  \end{enumerate}
\end{frame}
% ~~~~~~~~~~~~~~~~~~~~~~~~~~~~~~~~~~~~~~ A frm




% ~~~~~~~~~~~~~~~~~~~~~~~~~~~~~~~~~~~~~~ V frm
\begin{frame}[t,fragile]
	\begin{center}
		\vspace{4cm}
		\hRemark{\Huge Questions?}
	\end{center}
\end{frame}
% ~~~~~~~~~~~~~~~~~~~~~~~~~~~~~~~~~~~~~~ A frm



\end{document}
